\documentclass[10pt, letterpaper]{article}

\usepackage{CJKutf8}
\usepackage{resume}
\hypersetup{
    pdftitle={Jerome Gao - 全栈工程师},
    pdfsubject={全栈工程师简历},
    unicode=true
}

\begin{document}
\begin{CJK*}{UTF8}{gbsn}

    \newcommand{\AND}{\unskip\cleaders\copy\ANDbox\hskip\wd\ANDbox\ignorespaces}
    \newsavebox\ANDbox
    \sbox\ANDbox{$|$}

    \begin{header}
        \fontsize{25 pt}{25 pt}\selectfont Jerome Gao

        \vspace{5 pt}

        \normalsize
        \mbox{新西兰}%
        \kern 5.0 pt\AND\kern 5.0 pt
        \mbox{\href{mailto:gaojianqi6@gmail.com}{gaojianqi6@gmail.com}}%
        \kern 5.0 pt\AND\kern 5.0 pt
        \mbox{\href{tel:+64-027-493-0491}{+64 027 493 0491}}%
        \kern 5.0 pt\AND\kern 5.0 pt
        \mbox{\href{https://linkedin.com/in/jianqi6}{linkedin.com/in/jianqi6}}%
        \kern 5.0 pt\AND\kern 5.0 pt
        \mbox{\href{https://github.com/gaojianqi6}{github.com/gaojianqi6}}%
    \end{header}

    \vspace{5 pt - 0.3 cm}

    \section{职业概述}
        \begin{onecolentry}
            拥有十年以上经验的全栈工程师，能够端到端交付 Web 与跨端产品。后端覆盖 Node.js、Nest.js、Express、ASP.NET Core、Python、Java 和 PHP，数据库包括 PostgreSQL、MySQL 与 MongoDB；客户端使用 React、Next.js、Vue/Nuxt 和 React Native。熟悉 API 合约、鉴权、校验、异步任务、测试、容器化、云交付和运行可见性。新西兰毕业后工作签证有效期至 2028 年 12 月 02 日。
        \end{onecolentry}

    \section{核心技能}
        \begin{onecolentry}
            \textbf{后端：} Node.js、Nest.js、Express.js、ASP.NET Core、Python FastAPI、Java (Spring/Struts/Hibernate)、PHP、REST API、GraphQL、WebSocket
        \end{onecolentry}
        \vspace{0.15 cm}
        \begin{onecolentry}
            \textbf{数据库：} PostgreSQL、MySQL、MongoDB、Entity Framework Core、Prisma、Mongoose
        \end{onecolentry}
        \vspace{0.15 cm}
        \begin{onecolentry}
            \textbf{前端与移动端：} React、Next.js、TypeScript、JavaScript、Redux、Zustand、React Native、Expo、Vue/Nuxt、HTML5、CSS3
        \end{onecolentry}
        \vspace{0.15 cm}
        \begin{onecolentry}
            \textbf{测试：} Jest、React Testing Library、Playwright、Cypress、Vitest、xUnit、集成测试与 E2E 测试
        \end{onecolentry}
        \vspace{0.15 cm}
        \begin{onecolentry}
            \textbf{工具与 DevOps：} Docker、GitHub Actions、Azure、AWS、GCP、Vercel、Git、Lerna、Monorepo、结构化日志、可观测性
        \end{onecolentry}
        \vspace{0.15 cm}
        \begin{onecolentry}
            \textbf{AI 工程：} Gemini 集成、Schema 约束工作流、Claude Code、Codex、Cursor
        \end{onecolentry}

    \section{工作经历}
        \begin{twocolentry}{2025.07 -- 2025.11}
            \textbf{Full Stack Engineer}，Te Tāwharau o te Whakatōhea
        \end{twocolentry}
        \begin{onecolentry}
            \begin{highlights}
                \item 使用 React Native、Expo、Supabase 和 PostgreSQL 构建 Te Reo Māori 学习产品的棋盘游戏功能。
                \item 交付房间、实时回合、游戏状态、积分，以及跨设备的单人和多人模式。
            \end{highlights}
        \end{onecolentry}

        \vspace{0.25 cm}
        \begin{twocolentry}{2021.04 -- 2023.11}
            \textbf{Frontend Engineer / 全栈贡献}，Carsome
        \end{twocolentry}
        \begin{onecolentry}
            \begin{highlights}
                \item 交付 Nuxt SSR 官网、CMS 微前端、共享 TypeScript 包和覆盖车辆生命周期的运营流程。
                \item 开发 Nest.js 库存服务，使用 MongoDB、Redis 与 Bull 队列处理入库、中转、放行、维修、盘点、二维码操作和下游系统同步。
                \item 在 React、Vue 3、Angular.js 与既有 PHP CodeIgniter 系统之间协作，保持清晰职责和独立发布边界。
                \item 项目资料记录了 200+ CMS 日活用户，以及平台中管理的 10,000+ 辆车辆。
            \end{highlights}
        \end{onecolentry}

        \vspace{0.25 cm}
        \begin{twocolentry}{2020.07 -- 2020.12}
            \textbf{Software Engineer}，创量
        \end{twocolentry}
        \begin{onecolentry}
            \begin{highlights}
                \item 使用 Vue.js 2 与 Element UI 构建快手、抖音和百度广告运营流程。
                \item 通过共享表单模式、校验、Lint 和评审规范提升可维护性。
            \end{highlights}
        \end{onecolentry}

        \vspace{0.2 cm}
        \begin{twocolentry}{2017.06 -- 2020.02}
            \textbf{Software Engineer / 全栈贡献}，Pintec
        \end{twocolentry}
        \begin{onecolentry}
            \begin{highlights}
                \item 使用 Node.js、Express 与 MySQL 构建 API/BFF，覆盖鉴权、校验、KYC、风险测评、投资组合、订单和再平衡流程。
                \item 交付 React Web、移动 Web 与 React Native 客户端，通过共享 Redux Actions、Reducers、Selectors、API Model 和校验规则复用约 80\% 的业务逻辑。
                \item 接入专业投资算法和银行合作方合约，并保持明确的系统集成边界。
                \item 同期交付购它商城：面向响应式网站和移动 WebView 的 Node.js、Express、MySQL 与 React 电商产品。
            \end{highlights}
        \end{onecolentry}

        \vspace{0.2 cm}
        \begin{twocolentry}{2014.12 -- 2016.05}
            \textbf{Full Stack Engineer}，Cmku / 积木盒子
        \end{twocolentry}
        \begin{onecolentry}
            \begin{highlights}
                \item 构建连接 Web 客户端与 Java 服务的 Node.js/Express 集成层，以及 React/Redux 移动借贷流程和运营 CMS。
                \item 通过 UA 路由和服务端模板交付独立的移动端与桌面端体验。
            \end{highlights}
        \end{onecolentry}

        \vspace{0.2 cm}
        \begin{twocolentry}{2011 -- 2014}
            \textbf{Full Stack Engineer}，Cnlive / HelpNow / 六合
        \end{twocolentry}
        \begin{onecolentry}
            \begin{highlights}
                \item 使用 Spring、Struts、Hibernate 和 MySQL 开发 Java 全栈应用，包括实时看板与数据分析。
                \item 实现 HTML5 WebRTC 音视频通话，并通过 WebSocket 接入 C++ 信令服务。
            \end{highlights}
        \end{onecolentry}

    \section{独立产品}
        \begin{twocolentry}{2026 -- 至今}
            \textbf{Mealway -- AI 饮食计划产品}
        \end{twocolentry}
        \begin{onecolentry}
            \begin{highlights}
                \item 构建 ASP.NET Core 10 Minimal API、EF Core、PostgreSQL、Firebase JWT、Gemini 工作流和共享 Next.js/Expo 客户端合约。
                \item 实现模型选择、限流、Fallback、结构化日志、Docker、Azure Container Apps、Key Vault、Supabase、Vercel 和 xUnit 集成测试。
            \end{highlights}
        \end{onecolentry}

        \vspace{0.15 cm}
        \begin{twocolentry}{\href{https://rating.ospreypulse.com}{rating.ospreypulse.com}}
            \textbf{RateEverything -- 全栈内容评分平台}
        \end{twocolentry}
        \begin{onecolentry}
            \begin{highlights}
                \item 构建 Next.js 用户端与后台 CMS、Nest.js API、Python FastAPI 管理服务、共享 PostgreSQL 数据模型、测试和 CI/CD。
            \end{highlights}
        \end{onecolentry}

    \section{教育背景}
        \begin{twocolentry}{2024.07 -- 2025.11}
            \textbf{Master of Information Technology (Applied Computing)}，University of Waikato，新西兰
        \end{twocolentry}
        \vspace{0.10 cm}
        \begin{twocolentry}{2009.09 -- 2012.12}
            \textbf{信息系统工程学士}，国防科技大学，中国
        \end{twocolentry}

\end{CJK*}
\end{document}
